\section{Introduction}\label{ch5:sec:Introduction}

Energy efficiency has emerged as a fundamental constraint in modern high-performance computing (HPC) systems. Large production systems execute long sequences of computational jobs under strict power budgets, where electricity costs can dominate total cost of ownership and, in many cases, overtake hardware costs by a large margin. At the same time, HPC applications are typically subject to performance constraints such as deadlines, turnaround time requirements, or bounded performance degradation. These competing objectives make energy-aware runtime control a central challenge in today’s HPC environments.

A common abstraction used to study this problem is that of executing a sequence of jobs on a variable-speed processor so as to minimize energy consumption while satisfying performance constraints. In this setting, jobs often have strict deadlines but do not necessarily require a fixed execution time. Instead, applications can tolerate controlled slowdowns as long as they complete within user-specified or system-defined limits. Dynamic speed scaling, enabled by modern processors through mechanisms such as frequency control and power capping, provides a principled way to trade performance for energy savings by adjusting processor speed according to demand.

Dynamic speed scaling introduces a fundamental multi-objective optimization problem. Across different hardware platforms and applications, one constraint remains constant: energy consumption must be reduced. At the same time, performance requirements must be met. In current HPC systems, schedulers are responsible for deciding which resources to allocate to a job and, in some cases, selecting operating parameters such as node type or power limits. However, schedulers typically operate with limited knowledge of application-level performance sensitivity and must make decisions under uncertainty.

Much of the foundational work assumes an offline setting, where all jobs, along with their characteristics, are known in advance. This allows the computation of optimal or near-optimal schedules using global optimization techniques. In contrast, real HPC systems largely operate in an online setting, where jobs arrive dynamically and decisions must be made instantaneously without knowledge of future workloads. Online algorithms are commonly evaluated using competitive analysis, where an algorithm is said to be c-competitive if its objective value is within a factor c of the optimal offline solution for the same input.

To bridge theory and practice, HPC research often relies on proxy applications that capture representative computation and communication patterns of real workloads. Studying these proxy applications under different power caps allows us to understand how application performance degrades with reduced power and how these relationships vary across applications and hardware platforms. Such empirical characterization is essential for designing practical runtime controllers.

Reinforcement learning (RL) has gained attention as a promising framework for addressing this trade-off due to its ability to learn control policies directly from data. However, online RL is often impractical in HPC environments because exploration can degrade application performance and interfere with production workloads. Offline RL has therefore been explored as an alternative, where policies are trained using pre-collected system and application traces. Many existing approaches formulate the problem as a multi-objective optimization and convert it into a single objective by fixing weights on performance and energy, which often requires training multiple agents to convert different operating points.

In this paper, we focus on preference-driven control of HPC application runtime using a modified version of Decay-Driven Multi-Objective Reinforcement Learning (DDMORL) framework. Our goal is to provide a single, application- and hardware-agnostic control that allows users to explicitly specify their acceptable performance degradation through a simple preference input. By leveraging mathematical models of application behavior on HPC nodes to generate preference signals for augmenting the training data, we remove the need for interpolation models during training and enable efficient offline learning. This approach eliminates the requirement to train separate RL agents for different performance-energy trade-offs.

Through extensive experiments on multiple applications and compute nodes, we demonstrate that user-defined performance degradation levels can be reliably mapped to power caps that achieve high energy efficiency. The results show that the proposed methods delivers significant energy savings while maintaining performance within the user's specified limits, making it suitable for deployment in production HPC environments.